%%%%%%%%%%%%%%%%%%%%%%%%%%%%%%%%%%%%%%%%%%%%%%%%%%%%%%%%%%%%%%%%%%%%%%%%%%%%%%
%  Unit V, second half: Digital Analytics (Chapters 5-7)
%%%%%%%%%%%%%%%%%%%%%%%%%%%%%%%%%%%%%%%%%%%%%%%%%%%%%%%%%%%%%%%%%%%%%%%%%%%%%%

\chapter{The Evolution of Digital Analytics}

\syllabusstrip{Digital Analytics: Evolution of Digital Analytics.}

\begin{learningoutcomes}
\begin{itemize}
\item What digital analytics is, and how it differs from web analytics.
\item How digital analytics influences business decisions.
\item The six eras of the discipline --- what each measured, with what tools, and
      what each could not do.
\item The direction of travel across all six.
\end{itemize}
\end{learningoutcomes}

\section{Definition}

\begin{definitionbox}[Digital Analytics]
\keytermas{Digital analytics}{digital analytics} is the collection, measurement,
analysis and interpretation of digital data in order to improve business
performance and customer experience.
\pyq{CIE-III, Jun 2026 --- 2 M}
\end{definitionbox}

\begin{keybox}[Digital Analytics Is Not Web Analytics]
The distinction matters and is worth stating explicitly in any answer.

\smallskip
\textbf{Web analytics} is confined to website behaviour. \textbf{Digital
analytics} is broader --- it unifies website, mobile app, social, e-mail, paid
media, CRM, offline and support data into a single view of the customer, and its
purpose is to \emph{inform business decisions} rather than merely to report on a
channel.
\end{keybox}

\section{How Digital Analytics Influences Business Decisions}

\begin{exambox}[Model Answer --- 2 Marks]
\pyq{CIE-III, Jun 2026 --- 2 M}
Digital analytics provides customer insights and performance data that help
businesses improve marketing strategies and customer engagement.

\smallskip
Concretely, it determines where budget is allocated, which products are
developed, which segments are targeted, which experiences are fixed first, and
what the business forecasts.
\end{exambox}

\section{The Six Eras}

\begin{keybox}[High Yield]
``Trace the evolution of digital analytics and its growing importance in
delivering a seamless end-to-end customer experience.''
\pyq{Improvement CIE, 4 Jun 2025 --- 10 M}
Give the eras \emph{with their limitations}: the limitation of each era is what
motivated the next, and that causal chain is the argument the question is asking
for.
\end{keybox}

\begin{figure}[htbp]
\centering
\begin{tikzpicture}[font=\sffamily\footnotesize,
    era/.style={draw=ocre,line width=0.85pt,fill=ocrepale,rounded corners=3pt,
                minimum height=1.5cm,text width=1.55cm,align=center,
                font=\sffamily\bfseries\scriptsize,text=slate,inner sep=3pt}]
  \draw[rulegrey,line width=1pt] (-0.4,0) -- (11.6,0);

  \node[era] (e1) at (0.55,1.15)  {1\\ server logs\\[-1pt]{\scriptsize\mdseries 1990s}};
  \node[era] (e2) at (2.75,-1.15) {2\\ page-tag web analytics\\[-1pt]{\scriptsize\mdseries early 2000s}};
  \node[era] (e3) at (4.75,1.15)  {3\\ visitor \& goal analytics\\[-1pt]{\scriptsize\mdseries mid 2000s}};
  \node[era] (e4) at (6.75,-1.15) {4\\ multi-channel \& social\\[-1pt]{\scriptsize\mdseries 2010s}};
  \node[era] (e5) at (8.75,1.15)  {5\\ user-centric, event-based\\[-1pt]{\scriptsize\mdseries late 2010s}};
  \node[era,draw=greenok,fill=greenpale] (e6) at (10.85,-1.15)
     {6\\ unified, predictive, privacy-first\\[-1pt]{\scriptsize\mdseries present}};

  \foreach \n/\x in {e1/0.55,e3/4.75,e5/8.75}{
    \draw[ocre,line width=0.8pt] (\x,0) -- (\x,0.4);
    \fill[ocre] (\x,0) circle (0.075);}
  \foreach \n/\x in {e2/2.75,e4/6.75}{
    \draw[ocre,line width=0.8pt] (\x,0) -- (\x,-0.4);
    \fill[ocre] (\x,0) circle (0.075);}
  \draw[greenok,line width=0.8pt] (10.85,0) -- (10.85,-0.4);
  \fill[greenok] (10.85,0) circle (0.075);

  \node[anchor=north,text=slatelight,font=\sffamily\scriptsize\itshape,
        align=center,text width=11.8cm] at (5.6,-2.35)
    {counting files \ra{} counting visits \ra{} counting conversions \ra{}
     understanding people \ra{} predicting behaviour};
\end{tikzpicture}
\caption[The six eras of digital analytics]%
        {Six eras. Each was created by the limitation of the one before it, and
         each moved analytics closer to the customer and closer to the boardroom.}
\label{fig:sixeras}
\end{figure}

\begin{mmlongtable}{The six eras of digital analytics}{tab:sixeras}%
{@{}P{0.167\textwidth}P{0.245\textwidth}P{0.235\textwidth}P{0.265\textwidth}@{}}
\rowcolor{ocre}\thd{Era} & \thd{What was measured} & \thd{Tools and methods} & \thd{Limitation}\\
\midrule
\endfirsthead
\rowcolor{ocre}\thd{Era} & \thd{What was measured} & \thd{Tools and methods} & \thd{Limitation}\\
\midrule
\endhead
\rlab{1. Server log analysis} \newline {\scriptsize 1990s} & Hits, file requests,
bandwidth & Raw server logs, log parsers & Counted files, not people; caches and
bots corrupted the data\\
\rowcolor{tablealt}
\rlab{2. Page-tag web analytics} \newline {\scriptsize early 2000s} & Pageviews,
visits, referrers, bounce rate & JavaScript tags; the first hosted analytics
platforms & Session-based and site-siloed; no view of the person across visits\\
\rlab{3. Visitor and goal analytics} \newline {\scriptsize mid to late 2000s} &
Goals, funnels, e-commerce revenue, segments & Goal tracking, funnel
visualisation, A/B testing tools & Last-click attribution; ignored the rest of the
journey\\
\rowcolor{tablealt}
\rlab{4. Multi-channel and social analytics} \newline {\scriptsize 2010s} &
Cross-channel attribution, social listening, mobile app analytics & Multi-channel
funnels, mobile SDKs, social insight platforms & Data still fragmented across
separate platforms\\
\rlab{5. User-centric and event-based analytics} \newline {\scriptsize late 2010s
onward} & The individual user across devices and sessions; every interaction as an
event & Event-based analytics platforms, tag managers, cross-device identity &
Requires disciplined measurement planning and governance\\
\rowcolor{tablealt}
\rlab{6. Unified, predictive and privacy-first analytics} \newline {\scriptsize
present} & End-to-end customer experience; predicted churn and lifetime value;
consent-governed first-party data & Customer data platforms, data warehouses,
machine learning, server-side tagging, consent management & Privacy regulation and
cookie deprecation constrain tracking; modelling replaces some observation\\
\end{mmlongtable}

\begin{keybox}[The Direction of Travel]
The movement is consistent across all six eras: from \textbf{counting files
\ra{} counting visits \ra{} counting conversions \ra{} understanding people
\ra{} predicting behaviour}, and from \textbf{channel-siloed reporting \ra{} a
unified end-to-end view}.

\smallskip
Each shift moved analytics closer to the customer and closer to the boardroom.
\end{keybox}

\begin{summarybox}
Digital analytics is the collection, measurement, analysis and interpretation of
digital data to improve business performance and customer experience --- broader
than web analytics because it unifies web, app, social, e-mail, paid, CRM, offline
and support data, and because its purpose is decision rather than report. It
influences budget allocation, product development, segment targeting, experience
priorities and forecasting. Six eras mark its evolution: server logs, page-tag web
analytics, visitor and goal analytics, multi-channel and social analytics,
user-centric event-based analytics, and unified predictive privacy-first
analytics. Each era's limitation created the next, and the direction of travel runs
from counting files to predicting behaviour, and from silos to a unified view.
\end{summarybox}

%%%%%%%%%%%%%%%%%%%%%%%%%%%%%%%%%%%%%%%%%%%%%%%%%%%%%%%%%%%%%%%%%%%%%%%%%%%%%%
\chapter{The End-to-End Customer Experience}
%%%%%%%%%%%%%%%%%%%%%%%%%%%%%%%%%%%%%%%%%%%%%%%%%%%%%%%%%%%%%%%%%%%%%%%%%%%%%%

\syllabusstrip{Digital Analytics: information about end-to-end customer
experience.}

\begin{learningoutcomes}
\begin{itemize}
\item What the end-to-end customer experience means in digital analytics.
\item Why the siloed, channel-by-channel view fails.
\item The seven requirements for genuine end-to-end measurement.
\item The business value it produces.
\item How market segmentation is refined using digital analytics data.
\end{itemize}
\end{learningoutcomes}

\section{What It Means}

\begin{exambox}[Model Answer --- 2 Marks]
``What is meant by the `end-to-end customer experience' in digital analytics?''
\pyq{SEE, Jun/Jul 2025 --- 2 M}

The end-to-end customer experience means measuring and understanding the
\textbf{complete customer journey\ix{customer journey} across every touchpoint and every stage} ---
from the first anonymous impression\ix{impression}, through research, comparison, purchase,
delivery and support, to repeat purchase and advocacy\ix{advocacy} --- as a \textbf{single
connected experience} rather than as a set of disconnected channel reports.
\end{exambox}

\section{Why the Siloed View Fails}

\begin{itemize}
\item A visitor may discover the brand on one social platform, research it on
      video, compare it on search, read reviews on a marketplace and buy in a
      store. Channel-level reporting credits only the last touch, and therefore
      misleads budget allocation.
\item The modern journey is not a funnel. It is a \keytermas{decision
      constellation}{decision constellation} --- a network of micro-decisions
      happening in different places at the same time. Optimising for visibility in
      one place while the customer is deciding everywhere else produces the classic
      symptom: traffic looks healthy but conversions are flat.
\item Silos also hide experience failures that span teams --- for example, a
      marketing promise that the fulfilment process cannot keep.
\end{itemize}

\begin{figure}[htbp]
\centering
\begin{tikzpicture}[font=\sffamily\footnotesize]
  % --- left: the siloed view
  \node[anchor=north,text=warnred,font=\sffamily\bfseries\scriptsize] at (2.3,2.5)
    {THE SILOED VIEW};
  \foreach \i/\x in {1/0.6,2/1.75,3/2.9,4/4.05}{
    \fill[warnpale] (\x-0.42,-1.5) rectangle (\x+0.42,1.9);
    \draw[warnred!60,line width=0.8pt] (\x-0.42,-1.5) rectangle (\x+0.42,1.9);}
  \node[rotate=90,text=slate,font=\sffamily\scriptsize] at (0.6,0.2)  {search};
  \node[rotate=90,text=slate,font=\sffamily\scriptsize] at (1.75,0.2) {social};
  \node[rotate=90,text=slate,font=\sffamily\scriptsize] at (2.9,0.2)  {e-mail};
  \node[rotate=90,text=slate,font=\sffamily\scriptsize] at (4.05,0.2) {store};
  \node[anchor=north,text=slatelight,font=\sffamily\scriptsize\itshape,
        align=center,text width=4.4cm] at (2.3,-1.65)
    {four reports, four last-click winners, one customer nobody can see};

  % --- right: the unified view
  \node[anchor=north,text=greenok,font=\sffamily\bfseries\scriptsize] at (8.9,2.5)
    {THE UNIFIED VIEW};
  \draw[greenok,line width=1pt,rounded corners=5pt,fill=greenpale]
    (6.5,-1.5) rectangle (11.3,1.9);
  \foreach \a/\lbl in {150/search,90/social,30/video,-30/reviews,-90/store,-150/support}{
    \node[draw=greenok!70,line width=0.7pt,fill=white,rounded corners=2pt,
          inner sep=2.2pt,font=\sffamily\scriptsize,text=slate]
      (t\lbl) at ($(8.9,0.2)+(\a:1.72cm and 1.05cm)$) {\lbl};
    \draw[greenok!60,line width=0.6pt,dashed] (t\lbl) -- (8.9,0.2);}
  \node[circle,fill=greenok,text=white,inner sep=2.2pt,
        font=\sffamily\bfseries\scriptsize] at (8.9,0.2) {1};
  \node[anchor=north,text=slatelight,font=\sffamily\scriptsize\itshape,
        align=center,text width=4.8cm] at (8.9,-1.65)
    {one identity, one event taxonomy, multi-touch credit --- one journey};

  \draw[-{Latex[length=2.6mm]},slate,line width=1pt] (4.75,0.2) -- (6.35,0.2);
\end{tikzpicture}
\caption[The siloed view against the unified view]%
        {Four channel reports do not add up to one customer. The unified view is
         what makes multi-touch attribution, journey analysis and
         cross-channel personalisation possible at all.}
\label{fig:siloedunified}
\end{figure}

\section{What End-to-End Analytics Requires}

\begin{enumerate}
\item A \sfb{unified customer identity} across web, app, e-mail, CRM and offline
      systems.
\item A \sfb{consistent event taxonomy}, so that the same action means the same
      thing in every system.
\item \sfb{Journey and path analysis}, not just page reports.
\item \sfb{Multi-touch attribution} rather than last-click.
\item \sfb{Voice-of-customer data} --- surveys, net promoter score, reviews and
      support tickets --- joined to behavioural data.
\item \sfb{Post-purchase and service data}, so that the measurement does not stop
      at the transaction.
\item \sfb{Consent and privacy governance} built into the collection layer.
\end{enumerate}

\section{Business Value}

\begin{itemize}
\item \sfb{Accurate budget allocation}, based on true contribution rather than
      last-click credit.
\item \sfb{Identification of the highest-friction moment} in the whole journey,
      not merely on one page.
\item \sfb{Personalisation that is consistent} across channels rather than
      contradictory.
\item \sfb{Improved retention}, because the analysis continues past the purchase.
\item \sfb{Predictive capability} --- churn\ix{churn} risk and lifetime value modelled from
      journey behaviour.
\end{itemize}

\section{Refining Market Segmentation With Digital Analytics}

\begin{keybox}[High Yield]
``Analyze how market segmentation can be refined using digital analytics data.''
\pyq{SEE, Jun/Jul 2025, Q9b --- 8 M}
\end{keybox}

\begin{itemize}
\item \sfb{From demographic to behavioural segmentation} --- traditional
      segmentation groups by age and income; analytics groups by what people
      actually do, which predicts purchase far better.
\item \sfb{Value-based segmentation} --- \keytermas{RFM analysis}{RFM analysis}
      (recency, frequency, monetary value) and lifetime-value modelling identify
      the segments worth investing in.
\item \sfb{Journey-stage segmentation} --- separating first-time visitors,
      researchers, cart abandoners, first-time buyers and loyalists, each
      requiring different messaging.
\item \sfb{Intent-based segmentation} --- search query and on-site behaviour
      reveal purchase intent in real time.
\item \sfb{Channel-affinity segmentation} --- which segment responds to which
      channel, enabling efficient allocation.
\item \sfb{Predictive segmentation} --- machine learning models group users by
      predicted churn risk or predicted conversion probability rather than by past
      attributes.
\item \sfb{Continuous refinement} --- analytics-driven segments update
      automatically as behaviour changes, whereas survey-based segments decay from
      the day they are created.
\end{itemize}

\begin{summarybox}
The end-to-end customer experience means measuring the complete journey across
every touchpoint and stage as one connected experience rather than as disconnected
channel reports. The siloed view fails because last-click credit misleads budget
allocation, because the modern journey is a decision constellation rather than a
funnel, and because silos hide failures that span teams. Genuine end-to-end
measurement requires unified identity, a consistent event taxonomy, journey
analysis, multi-touch attribution, voice-of-customer data, post-purchase data and
privacy governance. Its value is accurate budget allocation, identification of the
worst friction in the whole journey, consistent personalisation, improved
retention and predictive capability. Analytics refines segmentation from
demographic to behavioural, value-based, journey-stage, intent-based,
channel-affinity and predictive --- and keeps those segments current
automatically.
\end{summarybox}

%%%%%%%%%%%%%%%%%%%%%%%%%%%%%%%%%%%%%%%%%%%%%%%%%%%%%%%%%%%%%%%%%%%%%%%%%%%%%%
\chapter{The Analyst's Influence and Role as a Change Agent}
%%%%%%%%%%%%%%%%%%%%%%%%%%%%%%%%%%%%%%%%%%%%%%%%%%%%%%%%%%%%%%%%%%%%%%%%%%%%%%

\syllabusstrip{Digital Analytics: analyst's influence on business, role as a
change agent.}

\begin{learningoutcomes}
\begin{itemize}
\item Why data requires interpretation, and where the chain can break.
\item The seven ways an analyst influences the business.
\item What separates a reporting function from a change agent.
\item What being a change agent requires, and the barriers it must overcome.
\item Designing a holistic marketing\ix{holistic marketing} dashboard.
\item Applying digital analytics to two worked business problems.
\end{itemize}
\end{learningoutcomes}

\section{Data Requires Interpretation}

A theme repeated throughout this unit: data requires interpretation, and errors
compound at every step of the chain.

\begin{figure}[htbp]
\centering
\begin{tikzpicture}[font=\sffamily\footnotesize,
    stp/.style={draw=slate!55,line width=0.85pt,fill=white,rounded corners=3pt,
                minimum height=1.05cm,text width=1.85cm,align=center,
                font=\sffamily\bfseries\scriptsize,text=slate,inner sep=3pt}]
  \node[stp] (c1) at (0,0)     {1. product design};
  \node[stp,fill=ocrepale,draw=ocre] (c2) at (2.35,0)  {2. tool setup};
  \node[stp] (c3) at (4.70,0)  {3. data output};
  \node[stp,fill=ocrepale,draw=ocre] (c4) at (7.05,0)  {4. human interpretation};
  \node[stp,fill=ocrepale,draw=ocre] (c5) at (9.40,0)  {5. business action};
  \foreach \a/\b in {c1/c2,c2/c3,c3/c4,c4/c5}{
    \draw[-{Latex[length=1.9mm]},slate!60,line width=0.85pt] (\a)--(\b);}

  \node[anchor=north,align=center,text width=2.15cm,font=\sffamily\scriptsize,
        text=warnred] at (0,-0.62)    {the product may be flawed};
  \node[anchor=north,align=center,text width=2.15cm,font=\sffamily\scriptsize,
        text=warnred] at (2.35,-0.62) {mis-tagged pages, duplicate tracking,
        wrong filters};
  \node[anchor=north,align=center,text width=2.15cm,font=\sffamily\scriptsize,
        text=warnred] at (4.70,-0.62) {sampling, bot traffic, consent gaps};
  \node[anchor=north,align=center,text width=2.15cm,font=\sffamily\scriptsize,
        text=warnred] at (7.05,-0.62) {correlation read as causation; averages
        hiding segments};
  \node[anchor=north,align=center,text width=2.15cm,font=\sffamily\scriptsize,
        text=warnred] at (9.40,-0.62) {the final decision may be wrong};

  \draw[decorate,decoration={brace,amplitude=5pt,mirror},ocre,line width=0.8pt]
    (1.35,-2.55) -- (10.4,-2.55);
  \node[anchor=north,text=ocredark,font=\sffamily\bfseries\scriptsize]
    at (5.9,-2.75) {the analyst's value lies in steps 2, 4 and 5};
\end{tikzpicture}
\caption[The interpretation chain]%
        {Five steps, each of which can break. The analyst does not control step~1
         and only partly controls step~3, which is why their contribution is
         concentrated in instrumenting, interpreting and translating.}
\label{fig:interpchain}
\end{figure}

\begin{mmlongtable}{Where the interpretation chain breaks}{tab:interpchain}%
{@{}P{0.19\textwidth}P{0.35\textwidth}P{0.39\textwidth}@{}}
\rowcolor{ocre}\thd{Step} & \thd{What happens} & \thd{Where it can go wrong}\\
\midrule
\endfirsthead
\rowcolor{ocre}\thd{Step} & \thd{What happens} & \thd{Where it can go wrong}\\
\midrule
\endhead
\rlab{1. Product design} & Someone designs the analytics tool & The product may be
flawed\\
\rowcolor{tablealt}
\rlab{2. Tool setup} & The tool is installed and preferences are set & The setup
may be wrong --- mis-tagged pages, duplicate tracking, wrong filters\\
\rlab{3. Data output} & The tool generates reports & Sampling, bot traffic,
consent gaps\\
\rowcolor{tablealt}
\rlab{4. Human interpretation} & Someone reads the data & The interpretation may
be wrong --- correlation read as causation, averages hiding segments\\
\rlab{5. Business action} & Reports are used to make choices & The final decision
may be wrong\\
\end{mmlongtable}

\begin{keybox}[The Analyst's Real Job]
A number is not a fact and a report is not an insight. The analyst's value lies in
steps 2, 4 and 5 --- instrumenting correctly, interpreting honestly, and
translating the interpretation into a decision the business can act on.

\smallskip
An organisation with excellent tools and poor interpretation makes confident
wrong decisions \emph{faster}.
\end{keybox}

\section{How the Analyst Influences the Business}

\begin{enumerate}
\item \sfb{Sets the measurement agenda.} By defining which KPIs are reported, the
      analyst determines what the organisation pays attention to.
\item \sfb{Arbitrates between opinions.} Analytics ends debates that seniority
      would otherwise decide.
\item \sfb{Reallocates money.} Channel and campaign analysis directly redirects
      marketing budget.
\item \sfb{Prioritises the roadmap.} Funnel analysis determines which product and
      site fixes are built first.
\item \sfb{Quantifies risk and opportunity.} Forecasts and cohort models inform
      pricing, inventory and hiring.
\item \sfb{Protects against self-deception.} The analyst is the person who reports
      that the favourite campaign did not work.
\item \sfb{Translates.} The analyst converts technical measurement into business
      language that executives can act on --- and this, more than technical skill,
      determines their influence.
\end{enumerate}

\section{The Analyst as a Change Agent}

\begin{keybox}[The Syllabus Names This Explicitly]
Reporting numbers is not change; changing what the organisation \emph{does} is.
\end{keybox}

\begin{mmlongtable}{From reporting function to change agent}{tab:changeagent}%
{@{}P{0.45\textwidth}P{0.48\textwidth}@{}}
\rowcolor{ocre}\thd{From (reporting function)} & \thd{To (change agent)}\\
\midrule
\endfirsthead
\rowcolor{ocre}\thd{From (reporting function)} & \thd{To (change agent)}\\
\midrule
\endhead
Produces scheduled reports on request & Proactively identifies problems nobody
asked about\\
\rowcolor{tablealt}
Answers ``what happened?'' & Answers ``why did it happen, and what should we
do?''\\
Describes performance & Recommends a specific action with an expected outcome\\
\rowcolor{tablealt}
Owned by one team & Embedded across marketing, product, sales and service\\
Measures after the fact & Designs experiments before decisions are made\\
\rowcolor{tablealt}
Presents data & Builds the business case and drives adoption of the decision\\
\end{mmlongtable}

\subsection{What being a change agent requires}

\begin{itemize}
\item \sfb{Business literacy} --- understanding the commercial model well enough
      to know which questions matter.
\item \sfb{Credibility} --- data quality and honest reporting, including
      inconvenient findings.
\item \sfb{Storytelling} --- a finding that cannot be communicated will not change
      anything.
\item \sfb{Stakeholder relationships} --- influence flows through trust, not
      through dashboards.
\item \sfb{A bias towards experimentation} --- proposing tests rather than
      asserting opinions.
\item \sfb{Follow-through} --- measuring whether the implemented change actually
      worked, and saying so.
\end{itemize}

\subsection{Barriers a change agent must overcome}

\begin{warnbox}[The Six Barriers]
\begin{itemize}
\item \sfb{Data distrust} --- ``the numbers must be wrong'' is the standard first
      response to unwelcome findings.
\item \sfb{HiPPO decision-making} --- the \keytermas{highest paid person's
      opinion}{HiPPO} overriding evidence.
\item \sfb{Vanity metrics} --- organisations preferring numbers that always go up.
\item \sfb{Siloed ownership} of data across departments.
\item \sfb{Privacy and consent constraints} limiting what can be measured.
\item \sfb{Analysis paralysis} --- endless measurement substituting for decision.
\end{itemize}
\end{warnbox}

\section{Designing a Holistic Marketing Dashboard}

\begin{keybox}[High Yield]
``Design a basic digital dashboard that reflects holistic marketing
performance.''
\pyq{SEE, Jun/Jul 2025, Q10a --- 8 M}
A good dashboard answers three questions in order --- \emph{are we growing?},
\emph{which channels are working?}, and \emph{where are we losing people?} --- and
every metric on it must be tied to a decision someone can make.
\end{keybox}

\begin{mmlongtable}{A holistic marketing dashboard}{tab:dashboard}%
{@{}P{0.20\textwidth}P{0.44\textwidth}P{0.29\textwidth}@{}}
\rowcolor{ocre}\thd{Dashboard section} & \thd{Metrics} & \thd{Decision it supports}\\
\midrule
\endfirsthead
\rowcolor{ocre}\thd{Dashboard section} & \thd{Metrics} & \thd{Decision it supports}\\
\midrule
\endhead
\rlab{1. Executive summary} & Revenue, conversions, conversion rate, customer
acquisition cost, return on advertising spend, lifetime value against acquisition
cost --- each against target and prior period & Is the business on plan?\\
\rowcolor{tablealt}
\rlab{2. Acquisition} & Sessions and conversions by channel; spend, cost per click
and cost per acquisition per channel; new against returning & Where should the
next rupee of budget go?\\
\rlab{3. Engagement} & Bounce rate, session duration, pages per session, scroll
depth, best and worst landing pages & Which content and pages need work?\\
\rowcolor{tablealt}
\rlab{4. Conversion funnel} & Step-by-step drop-off from landing to purchase,
split by device & Which single step do we fix first?\\
\rlab{5. Retention and loyalty} & Repeat purchase rate, cohort retention curves,
churn, net promoter score, review sentiment & Are we keeping the customers we
buy?\\
\rowcolor{tablealt}
\rlab{6. Channel deep-dives} & SEO --- impressions, rankings, organic sessions;
e-mail --- open rate, click-through rate, revenue per message; social --- reach,
engagement rate, share of voice; mobile --- installs, active users, retention &
Which channel-level tactic needs attention?\\
\rlab{7. Technical health} & Core Web Vitals, error rate, uptime, tracking
integrity & Is anything silently broken?\\
\end{mmlongtable}

\begin{keybox}[Dashboard Design Principles]
One screen, top-down importance. Show \textbf{trend}, not just the current value
--- a number without a comparison is meaningless. Always include the
\textbf{target} alongside the actual. \textbf{Segment} by device and channel,
because averages hide failures. Use \textbf{colour only to signal status}, never
for decoration. And apply the same rule as web design: \textbf{if a metric drives
no decision, delete it.}
\end{keybox}

\section{Applying Digital Analytics --- Two Worked Scenarios}

\begin{exambox}[Improving Customer Retention for a Food Delivery App]
\pyq{CIE-III, Jun 2026 --- 10 M}
\textbf{Official scheme answer.} The app can analyse customer ordering patterns,
feedback and abandoned carts. Personalised offers, loyalty programmes and targeted
notifications can improve retention. Analytics helps identify customer needs and
improve services effectively.

\smallskip
\textbf{A fuller answer, in five steps.}
\begin{enumerate}
\item \sfb{Measure} --- build cohort retention curves by acquisition source, and
      calculate churn rate, order frequency and average order value per segment.
\item \sfb{Segment} --- identify the behaviours that distinguish retained from
      churned users, such as ordering within the first seven days or using a saved
      address.
\item \sfb{Diagnose} --- analyse cart abandonment reasons, delivery-time
      complaints and support tickets to locate the experience failure.
\item \sfb{Act} --- trigger personalised offers based on past cuisine preference;
      run a win-back push after 14 inactive days; introduce a subscription pass for
      free delivery; add a loyalty streak to build habit; and fix the delivery-time
      promise that the data shows is most often broken.
\item \sfb{Measure again} --- track Day-30 retention, repeat order rate and
      lifetime value against a control group, to prove the change worked.
\end{enumerate}
\end{exambox}

\begin{exambox}[Maintaining Competitive Advantage in a Global Market]
\pyq{CIE-III, Jun 2026 --- 10 M}
\textbf{Official scheme answer.} Digital analytics helps understand customer
preferences, competitor trends and campaign effectiveness. Website optimization
improves user experience and conversions. Together, they support better
decision-making, customer engagement and global competitiveness.

\smallskip
\textbf{A fuller answer.} Analytics provides market-by-market insight into what
customers actually want, which channels perform in each region, and where
competitors are gaining share. Website optimization ensures that the traffic won
in each market \emph{converts} --- through localisation of language, currency,
payment methods and page speed on regional networks. The combination creates a
continuous feedback loop in which the business learns faster than its competitors.

\smallskip
\textbf{Speed of learning, not size, is the durable advantage in a global digital
market.}
\end{exambox}

\begin{summarybox}
Data requires interpretation, and the chain from product design through tool
setup, data output and human interpretation to business action can break at every
step; the analyst's value is concentrated in instrumenting, interpreting and
translating. The analyst influences the business by setting the measurement
agenda, arbitrating between opinions, reallocating money, prioritising the
roadmap, quantifying risk, protecting against self-deception and translating
measurement into business language. Becoming a change agent means moving from
scheduled reports to proactive problem-finding, from ``what happened'' to ``what
should we do'', and from presenting data to driving adoption --- which requires
business literacy, credibility, storytelling, relationships, a bias towards
experimentation and follow-through, against barriers of data distrust,
HiPPO decision-making, vanity metrics, silos, privacy constraints and analysis
paralysis. A holistic dashboard has seven sections, each tied to a decision, and
any metric that drives no decision should be deleted.
\end{summarybox}
